% =========================================================================
\section{Ejemplos}\label{sec:ejemplos}
% =========================================================================

Aplicamos ahora el método a cuatro ecuaciones de complejidad creciente. En
todas ellas $x_0=0$, de modo que $z=x$. La \cref{fig:saltos} anticipa la
estructura de las recurrencias que aparecerán: cuando $p$ y $q$ son
polinomios sencillos, la recurrencia «salta» de dos en dos o de tres en
tres y desacopla los coeficientes en cadenas independientes.

% Figura: cadenas de coeficientes desacopladas (saltos de 2 y de 3)
\begin{figure}[ht]
\centering
\begin{tikzpicture}
  \tikzset{coef/.style={circle, draw=#1, fill=cajafondo, line width=0.8pt,
                        minimum size=0.82cm, inner sep=0pt, font=\small},
           coef/.default=grismedio,
           cadena/.style={-{Stealth[length=5pt,width=4.5pt]}, line width=0.7pt, #1}}
  % ---------------- (a) Salto de dos ----------------
  \begin{scope}
    \node[font=\small, text=texto, anchor=west] at (-0.5,1.45) {(a) Salto de dos: $a_{n+2}\propto a_n$};
    \foreach \i/\k in {0/0,1/2,2/4,3/6}{
      \node[coef=cian] (p\i) at (1.35*\i,0.45) {$a_{\k}$};}
    \foreach \i/\k in {0/1,1/3,2/5,3/7}{
      \node[coef=magenta] (q\i) at (1.35*\i,-0.65) {$a_{\k}$};}
    \foreach \i in {0,1,2}{
      \pgfmathtruncatemacro{\j}{\i+1}
      \draw[cadena=cian] (p\i) -- (p\j);
      \draw[cadena=magenta] (q\i) -- (q\j);}
    \node[font=\small, text=gris] at (1.35*4-0.25,0.45) {$\cdots$};
    \node[font=\small, text=gris] at (1.35*4-0.25,-0.65) {$\cdots$};
    \node[etiqgris, anchor=north west, align=left] at (-0.5,-1.3)
      {cadena par $\propto a_0$, cadena impar $\propto a_1$\\
       $y''+y=0$, \ Legendre, \ $y''+xy'+y=0$};
  \end{scope}
  % ---------------- (b) Salto de tres ----------------
  \begin{scope}[xshift=7.6cm]
    \node[font=\small, text=texto, anchor=west] at (-0.5,1.45) {(b) Salto de tres: $a_{n+2}\propto a_{n-1}$};
    \foreach \i/\k in {0/0,1/3,2/6,3/9}{
      \node[coef=cian] (r\i) at (1.35*\i,0.45) {$a_{\k}$};}
    \foreach \i/\k in {0/1,1/4,2/7,3/10}{
      \node[coef=magenta] (s\i) at (1.35*\i,-0.65) {$a_{\k}$};}
    \foreach \i/\k in {0/2,1/5,2/8,3/11}{
      \node[coef, text=gris] (t\i) at (1.35*\i,-1.75) {$a_{\k}$};}
    \foreach \i in {0,1,2}{
      \pgfmathtruncatemacro{\j}{\i+1}
      \draw[cadena=cian] (r\i) -- (r\j);
      \draw[cadena=magenta] (s\i) -- (s\j);
      \draw[cadena=grismedio, densely dashed] (t\i) -- (t\j);}
    \foreach \y in {0.45,-0.65,-1.75}{
      \node[font=\small, text=gris] at (1.35*4-0.25,\y) {$\cdots$};}
    \node[etiqgris, anchor=west] at (1.35*4+0.15,-1.75) {$=0$};
    \node[etiqgris, anchor=north west, align=left] at (-0.5,-2.4)
      {$a_2=0$ anula la tercera cadena; ecuación de Airy};
  \end{scope}
\end{tikzpicture}
\caption{Cuando $p$ y $q$ son polinomios sencillos, la recurrencia
desacopla los coeficientes en cadenas independientes. (a) Si sólo
intervienen $a_n$ y $a_{n+2}$, la cadena par depende de $a_0$ y la impar de
$a_1$. (b) En la ecuación de Airy intervienen $a_{n-1}$ y $a_{n+2}$: hay tres
cadenas, y la que arranca en $a_2=0$ es idénticamente nula.}
\label{fig:saltos}
\end{figure}


% -------------------------------------------------------------------------
\subsection{Comprobación: el oscilador armónico}\label{sec:oscilador}
% -------------------------------------------------------------------------

\begin{ejemplo}[$y''+y=0$]\label{ej:oscilador}
Aquí $p\equiv0$ y $q\equiv1$, es decir, $q_0=1$ y el resto de coeficientes
son nulos. En \eqref{eq:recurrencia} sólo sobrevive el término $k=0$ de
$q$:
\begin{equation}\label{eq:rec-oscilador}
  (n+2)(n+1)\,a_{n+2}+a_n=0
  \quad\Longrightarrow\quad
  a_{n+2}=-\frac{a_n}{(n+2)(n+1)} .
\end{equation}
La recurrencia salta de dos en dos (\cref{fig:saltos}\,(a)) y una inducción
inmediata da $a_{2m}=(-1)^{m}a_0/(2m)!$ y $a_{2m+1}=(-1)^{m}a_1/(2m+1)!$:
\begin{itemize}
  \item con $a_0=1$, $a_1=0$: $a_2=-\frac1{2!}$, $a_4=\frac1{4!},\dots$, e
        impares nulos;
  \item con $a_0=0$, $a_1=1$: $a_3=-\frac1{3!}$, $a_5=\frac1{5!},\dots$, y
        pares nulos.
\end{itemize}
Por tanto
\begin{equation}\label{eq:cos-sen}
  y_1=\sum_{m=0}^{\infty}\frac{(-1)^{m}x^{2m}}{(2m)!}=\cos x,
  \qquad
  y_2=\sum_{m=0}^{\infty}\frac{(-1)^{m}x^{2m+1}}{(2m+1)!}=\sin x .
\end{equation}
Recuperamos lo que ya sabíamos, lo cual es una buena señal. Como $P\equiv1$
no tiene ceros, el radio es $R=\infty$ (\cref{fig:coseno}).
\end{ejemplo}

% Figura: sumas parciales de la serie del coseno (R = infinito)
\begin{figure}[ht]
\centering
\begin{tikzpicture}
  \begin{axis}[oscuro, width=0.92\textwidth, height=6.0cm,
      xmin=-7, xmax=7, ymin=-2.2, ymax=2.2,
      xtick={-6.2832,-3.1416,0,3.1416,6.2832},
      xticklabels={$-2\pi$,$-\pi$,$0$,$\pi$,$2\pi$},
      ytick={-2,-1,0,1,2}, xlabel={$x$},
      samples=400, domain=-7:7, restrict y to domain=-4:4,
      legend style={at={(0.5,-0.22)}, anchor=north, legend columns=5,
                    /tikz/every even column/.append style={column sep=8pt}}]
    \addplot[c2] {1-x^2/2};
    \addplot[c3] {1-x^2/2+x^4/24};
    \addplot[c4] {1-x^2/2+x^4/24-x^6/720+x^8/40320};
    \addplot[c5] {1-x^2/2+x^4/24-x^6/720+x^8/40320-x^10/3628800+x^12/479001600};
    \addplot[blanco, line width=1.6pt] {cos(deg(x))};
    \legend{$S_2$, $S_4$, $S_8$, $S_{12}$, $\cos x$}
  \end{axis}
\end{tikzpicture}
\caption{Sumas parciales $S_{2m}(x)=\sum_{k=0}^{m}(-1)^k x^{2k}/(2k)!$ de la
solución $y_1=\cos x$ del \cref{ej:oscilador}. Como $P\equiv1$ no tiene
ceros, el radio es infinito: al aumentar el orden, la aproximación es buena
en intervalos cada vez mayores, sin barrera alguna.}
\label{fig:coseno}
\end{figure}


% -------------------------------------------------------------------------
\subsection{La ecuación de Airy}\label{sec:airy}
% -------------------------------------------------------------------------

\begin{ejemplo}[$y''-x\,y=0$]\label{ej:airy}
Ahora $p\equiv0$ y $q=-x$, es decir, $q_1=-1$ y el resto nulos. En el
sumatorio de \eqref{eq:recurrencia} sólo sobrevive $k=1$, que existe si
$n\ge1$:
\begin{equation}\label{eq:rec-airy}
  (n+2)(n+1)\,a_{n+2}-a_{n-1}=0\quad(n\ge1),\qquad a_2=0\quad(n=0).
\end{equation}
La recurrencia salta de tres en tres, $a_{n+2}=a_{n-1}/\bigl((n+2)(n+1)\bigr)$
(\cref{fig:saltos}\,(b)):
\begin{equation}\label{eq:coef-airy}
  a_3=\frac{a_0}{3\cdot2},\quad a_6=\frac{a_0}{6\cdot5\cdot3\cdot2},\qquad
  a_4=\frac{a_1}{4\cdot3},\quad a_7=\frac{a_1}{7\cdot6\cdot4\cdot3},\qquad
  a_2=a_5=a_8=\dots=0 .
\end{equation}
En general, como $(3j-1)(3j)=9j\,(j-\tfrac13)$ y $3j\,(3j+1)=9j\,(j+\tfrac13)$,
\begin{equation}\label{eq:coef-airy-gamma}
  a_{3m}=\frac{\Gamma(\tfrac23)}{9^{m}\,m!\,\Gamma(m+\tfrac23)}\,a_0,
  \qquad
  a_{3m+1}=\frac{\Gamma(\tfrac43)}{9^{m}\,m!\,\Gamma(m+\tfrac43)}\,a_1 ,
\end{equation}
y el sistema fundamental es
\begin{equation}\label{eq:airy-y1y2}
  y_1=1+\frac{x^{3}}{6}+\frac{x^{6}}{180}+\cdots,\qquad
  y_2=x+\frac{x^{4}}{12}+\frac{x^{7}}{504}+\cdots ,
\end{equation}
con radio $R=\infty$. Estas funciones \emph{no son elementales}: el método de
series \emph{define} funciones nuevas. Las funciones de Airy de uso habitual
son las combinaciones \cite[§9.4]{DLMF}
\begin{equation}\label{eq:Ai-Bi}
  \Ai(x)=c_1\,y_1(x)-c_2\,y_2(x),\qquad
  \Bi(x)=\sqrt3\,\bigl(c_1\,y_1(x)+c_2\,y_2(x)\bigr),
\end{equation}
con $c_1=\Ai(0)=3^{-2/3}/\Gamma(\tfrac23)\approx0{,}355028$ y
$c_2=-\Ai'(0)=3^{-1/3}/\Gamma(\tfrac13)\approx0{,}258819$.
\end{ejemplo}

\begin{observacion}[Interpretación cualitativa]\label{obs:airy}
Para $x<0$ la ecuación es $y''=-\abs{x}\,y$: un oscilador de «frecuencia»
$\sqrt{\abs{x}}$ creciente, y las soluciones oscilan cada vez más deprisa
con amplitud decreciente. Para $x>0$ es $y''=x\,y$ y las soluciones crecen o
decrecen exponencialmente; $\Ai$ es la única (salvo múltiplos) que decae.
La ecuación de Airy es la ecuación de Schrödinger en un potencial lineal y
describe el comportamiento de las funciones de onda cerca de un punto de
retroceso clásico \cite[§10.2]{BENDER}.
\end{observacion}

% Figura: soluciones de la ecuación de Airy (datos: scripts/generar_datos.py)
\begin{figure}[ht]
\centering
\begin{tikzpicture}
  \begin{groupplot}[group style={group size=1 by 2, vertical sep=1.5cm},
      oscuro, width=0.92\textwidth, height=5.3cm,
      xmin=-9, xmax=2.4, xtick={-8,-6,-4,-2,0,2}, unbounded coords=jump]
    % ---------- (a) sistema fundamental ----------
    \nextgroupplot[ymin=-1.6, ymax=3.2, ytick={-1,0,1,2,3},
        title={(a) Sistema fundamental $y_1$, $y_2$ (series \eqref{eq:airy-y1y2})},
        legend pos=north west]
      \fill[grisoscuro, opacity=0.45] (axis cs:0,-1.6) rectangle (axis cs:2.4,3.2);
      \addplot[c1, restrict y to domain=-10:10] table[x=x, y=y1] {datos/airy.dat};
      \addplot[c2, restrict y to domain=-10:10] table[x=x, y=y2] {datos/airy.dat};
      \legend{$y_1$, $y_2$}
      \node[font=\scriptsize, text=gris, anchor=south] at (axis cs:-6.6,-1.58) {región oscilatoria ($x<0$)};
      \node[font=\scriptsize, text=gris, anchor=south] at (axis cs:1.2,-1.55) {crecimiento};
    % ---------- (b) funciones de Airy ----------
    \nextgroupplot[ymin=-0.65, ymax=1.6, ytick={-0.5,0,0.5,1,1.5},
        xlabel={$x$},
        title={(b) Funciones de Airy $\Ai$ y $\Bi$, combinaciones \eqref{eq:Ai-Bi}},
        legend pos=north west]
      \fill[grisoscuro, opacity=0.45] (axis cs:0,-0.65) rectangle (axis cs:2.4,1.6);
      \addplot[c3, restrict y to domain=-10:10] table[x=x,
          y expr={0.355028054*\thisrow{y1}-0.258819404*\thisrow{y2}}] {datos/airy.dat};
      \addplot[c4, restrict y to domain=-10:10] table[x=x,
          y expr={1.732050808*(0.355028054*\thisrow{y1}+0.258819404*\thisrow{y2})}] {datos/airy.dat};
      \legend{$\Ai$, $\Bi$}
  \end{groupplot}
\end{tikzpicture}
\caption{Ecuación de Airy $y''=x\,y$. Para $x<0$ las soluciones oscilan con
frecuencia creciente y amplitud decreciente; para $x>0$ (zona sombreada)
crecen exponencialmente, salvo $\Ai$, que decae. Las curvas se han obtenido
sumando directamente las series \eqref{eq:airy-y1y2}, que convergen en todo
$\R$.}
\label{fig:airy}
\end{figure}


% -------------------------------------------------------------------------
\subsection{Un problema de valores iniciales}\label{sec:pvi}
% -------------------------------------------------------------------------

\begin{ejemplo}[$y''+x\,y'+y=0$, $y(0)=y_0$, $y'(0)=v_0$]\label{ej:pvi}
Este problema procede de \cite{TEMAV}. Ahora $p=x$ y $q\equiv1$, es decir,
$p_1=1$ y $q_0=1$. En \eqref{eq:recurrencia} sobreviven el término $k=1$ de
$p$, que vale $n\,p_1\,a_n$, y el término $k=0$ de $q$:
\begin{equation}\label{eq:rec-pvi}
  (n+2)(n+1)\,a_{n+2}+(n+1)\,a_n=0
  \quad\Longrightarrow\quad
  a_{n+2}=-\frac{a_n}{n+2},\qquad n\ge0 .
\end{equation}
De nuevo la recurrencia salta de dos en dos. Con la notación del doble
factorial, $(2i)!!=2\cdot4\cdots(2i)=2^{i}\,i!$ y
$(2i+1)!!=1\cdot3\cdots(2i+1)$,
\begin{equation}\label{eq:coef-pvi}
  a_{2i}=\frac{(-1)^{i}}{(2i)!!}\,a_0,\qquad
  a_{2i+1}=\frac{(-1)^{i}}{(2i+1)!!}\,a_1 ,
\end{equation}
de donde
\begin{equation}\label{eq:pvi-y1y2}
  y_1=\sum_{i=0}^{\infty}\frac{(-1)^{i}}{(2i)!!}\,x^{2i}
     =\sum_{i=0}^{\infty}\frac{1}{i!}\Bigl(-\frac{x^{2}}{2}\Bigr)^{i}
     =\eul^{-x^{2}/2},
  \qquad
  y_2=\sum_{i=0}^{\infty}\frac{(-1)^{i}}{(2i+1)!!}\,x^{2i+1}.
\end{equation}
Imponiendo las condiciones iniciales, $a_0=y_0$ y $a_1=v_0$, y la solución
del problema es $y=y_0\,y_1+v_0\,y_2$, con radio $R=\infty$.
\end{ejemplo}

\begin{observacion}[Comprobación por integración directa]\label{obs:pvi-exacta}
La ecuación es exacta: $y''+x\,y'+y=\bigl(y'+x\,y\bigr)'$. Integrando una vez
con los datos iniciales, $y'+x\,y=v_0$, ecuación lineal de primer orden con
factor integrante $\eul^{x^{2}/2}$:
\[
  y(x)=\eul^{-x^{2}/2}\Bigl(y_0+v_0\int_0^{x}\eul^{t^{2}/2}\,\dif t\Bigr).
\]
Por la unicidad del \cref{teo:eyu}, $y_1=\eul^{-x^{2}/2}$ (confirmando
\eqref{eq:pvi-y1y2}) e
$y_2=\eul^{-x^{2}/2}\int_0^{x}\eul^{t^{2}/2}\,\dif t$, que es, salvo
reescalado, la función de Dawson: $y_2(x)=\sqrt2\,F(x/\sqrt2)$ con
$F(u)=\eul^{-u^{2}}\int_0^{u}\eul^{s^{2}}\dif s$ \cite[§7.2]{DLMF}. El método
de series ha reconstruido, sin integrar, una gaussiana y una función no
elemental.
\end{observacion}

% Figura: soluciones de y'' + x y' + y = 0 (datos: scripts/generar_datos.py)
\begin{figure}[ht]
\centering
\begin{tikzpicture}
  \begin{axis}[oscuro, width=0.92\textwidth, height=5.6cm,
      xmin=-5, xmax=5, ymin=-1.15, ymax=1.65,
      xtick={-4,-2,0,2,4}, ytick={-1,-0.5,0,0.5,1,1.5}, xlabel={$x$},
      legend style={at={(0.985,0.04)}, anchor=south east}]
    \addplot[c1] table[x=x, y=y1] {datos/ejemplo_gauss.dat};
    \addplot[c2] table[x=x, y=y2] {datos/ejemplo_gauss.dat};
    \addplot[c3, densely dashed] table[x=x, y expr={\thisrow{y1}+\thisrow{y2}}] {datos/ejemplo_gauss.dat};
    \legend{$y_1=\eul^{-x^2/2}$, $y_2=\eul^{-x^2/2}\!\int_0^x\eul^{t^2/2}\dif t$,
            $y_1+y_2$ \ ($y_0=v_0=1$)}
  \end{axis}
\end{tikzpicture}
\caption{Sistema fundamental del \cref{ej:pvi}: la gaussiana $y_1$ (par) y
la función de tipo Dawson $y_2$ (impar, con $y_2\sim1/x$ para $\abs{x}\to\infty$),
junto con la solución del problema de valores iniciales con $y_0=v_0=1$.}
\label{fig:gauss}
\end{figure}


% -------------------------------------------------------------------------
\subsection{La ecuación de Legendre}\label{sec:legendre}
% -------------------------------------------------------------------------
\paragraph{Truco práctico: no dividir por $P$}
Si $P$ es un polinomio, suele ser más cómodo sustituir la serie
directamente en \eqref{eq:edo}, sin pasar a la forma normal: así se evita
desarrollar los cocientes $A/P$ y $B/P$. Las multiplicaciones por potencias
de $x$ sólo desplazan índices: por ejemplo,
$x^{2}\,y''=\sum_{n\ge0}n(n-1)\,a_n\,x^{n}$ y $x\,y'=\sum_{n\ge0}n\,a_n\,x^{n}$
(los términos añadidos con $n=0,1$ son nulos).

\begin{ejemplo}[Legendre en $x_0=0$]\label{ej:legendre}
Consideremos
\begin{equation}\label{eq:legendre}
  (1-x^{2})\,y''-2x\,y'+\ell(\ell+1)\,y=0 .
\end{equation}
Sustituyendo y agrupando potencias $x^{n}$ ---el término $-x^{2}y''$ aporta
$-n(n-1)\,a_n$ y $-2x\,y'$ aporta $-2n\,a_n$---,
\begin{equation}\label{eq:rec-legendre-0}
  (n+2)(n+1)\,a_{n+2}-\bigl[n(n+1)-\ell(\ell+1)\bigr]\,a_n=0 ,
\end{equation}
y, factorizando $n(n+1)-\ell(\ell+1)=(n-\ell)(n+\ell+1)$,
\begin{equation}\label{eq:rec-legendre}
  \boxed{\;a_{n+2}=\frac{(n-\ell)(n+\ell+1)}{(n+2)(n+1)}\,a_n ,\qquad n\ge0.\;}
\end{equation}
El radio de convergencia garantizado es $R=1$, distancia de $0$ a los ceros
$x=\pm1$ de $P$. Para una serie que no se corte es exactamente $1$: como
$a_{n+2}/a_n\to1$, el criterio del cociente aplicado a la serie en $x^{2}$
da divergencia para $\abs{x}>1$.
\end{ejemplo}

\begin{proposicion}[Polinomios de Legendre]\label{prop:polinomios-legendre}
Si $\ell\in\N$, la serie de la misma paridad que $\ell$ se trunca: el
numerador de \eqref{eq:rec-legendre} se anula para $n=\ell$, luego
$a_{\ell+2}=a_{\ell+4}=\dots=0$, y la solución correspondiente es un
polinomio de grado $\ell$. Normalizado por $P_\ell(1)=1$, es el
\emph{polinomio de Legendre} $P_\ell$:
\[
  P_0=1,\quad P_1=x,\quad P_2=\tfrac12(3x^{2}-1),\quad
  P_3=\tfrac12(5x^{3}-3x),\quad P_4=\tfrac18(35x^{4}-30x^{2}+3).
\]
La serie de la otra paridad no se corta y tiene radio exactamente $1$.
\end{proposicion}

\begin{proof}
El truncamiento es inmediato a partir de \eqref{eq:rec-legendre}. Por
ejemplo, para $\ell=4$: $a_2=\frac{(-4)\cdot5}{2}\,a_0=-10\,a_0$,
$a_4=\frac{(-2)\cdot7}{12}\,a_2=\frac{35}{3}\,a_0$ y $a_6=0$; imponiendo
$a_0\bigl(1-10+\frac{35}3\bigr)=1$ se obtiene $a_0=\frac38$ y la expresión
de $P_4$. Si la otra serie se cortara, ambas soluciones serían polinomios y
estarían acotadas en $[-1,1]$; el apartado siguiente muestra (para $\ell=1$)
que no es así, y el caso general se trata en \cite[§15.6]{ARFKEN}.
\end{proof}

% Figura: polinomios de Legendre y serie no truncada para l = 1
\begin{figure}[ht]
\centering
\begin{tikzpicture}
  \begin{groupplot}[group style={group size=2 by 1, horizontal sep=1.3cm},
      oscuro, width=0.5\textwidth, height=6.2cm, xmin=-1, xmax=1,
      xtick={-1,-0.5,0,0.5,1}, xlabel={$x$}]
    % ---------- (a) P_0 ... P_4 ----------
    \nextgroupplot[ymin=-1.1, ymax=1.1, ytick={-1,-0.5,0,0.5,1},
        title={(a) Polinomios $P_\ell$ (series truncadas)},
        samples=200, domain=-1:1,
        legend style={at={(0.5,-0.24)}, anchor=north, legend columns=5,
                      /tikz/every even column/.append style={column sep=3pt}}]
      \addplot[line width=1.15pt, gris] {1};
      \addplot[c1] {x};
      \addplot[c2] {0.5*(3*x^2-1)};
      \addplot[c3] {0.5*(5*x^3-3*x)};
      \addplot[c4] {(35*x^4-30*x^2+3)/8};
      \legend{$P_0$, $P_1$, $P_2$, $P_3$, $P_4$}
    % ---------- (b) y1 = 1 - x artanh x ----------
    \nextgroupplot[ymin=-2.6, ymax=1.25, ytick={-2,-1,0,1},
        title={(b) Serie par para $\ell=1$: no se corta},
        legend style={at={(0.5,-0.24)}, anchor=north, legend columns=5,
                      /tikz/every even column/.append style={column sep=3pt}}]
      \draw[guia] (axis cs:-1,-2.6) -- (axis cs:-1,1.25);
      \draw[guia] (axis cs:1,-2.6) -- (axis cs:1,1.25);
      \addplot[c2] table[x=x, y=S1] {datos/legendre_l1_parciales.dat};
      \addplot[c3] table[x=x, y=S2] {datos/legendre_l1_parciales.dat};
      \addplot[c4] table[x=x, y=S4] {datos/legendre_l1_parciales.dat};
      \addplot[c5] table[x=x, y=S8] {datos/legendre_l1_parciales.dat};
      \addplot[blanco, line width=1.6pt, samples=400, domain=-0.997:0.997]
        {1-0.5*x*ln((1+x)/(1-x))};
      \legend{$S_1$, $S_2$, $S_4$, $S_8$, $y_1$}
  \end{groupplot}
\end{tikzpicture}
\caption{(a) Para $\ell\in\N$ la serie de la paridad de $\ell$ se corta y da
el polinomio de Legendre $P_\ell$, regular en todo $[-1,1]$ (y, de hecho, en
todo $\R$). (b) Para $\ell=1$, la serie par
$y_1=1-x\artanh x$ de \eqref{eq:legendre-q1} y sus sumas parciales
$S_M=1-\sum_{k=1}^{M}x^{2k}/(2k-1)$: la convergencia se degrada al
acercarse a $x=\pm1$, donde $y_1$ diverge logarítmicamente. El radio $R=1$
del \cref{teo:fuchs} se alcanza.}
\label{fig:legendre}
\end{figure}


\begin{ejemplo}[Legendre con $\ell=1$]\label{ej:legendre-1}
Para $(1-x^{2})\,y''-2x\,y'+2y=0$ \cite{TEMAV}, la recurrencia
\eqref{eq:rec-legendre} se reduce a
\begin{equation}\label{eq:rec-legendre-1}
  a_{n+2}=\frac{(n-1)(n+2)}{(n+2)(n+1)}\,a_n=\frac{n-1}{n+1}\,a_n .
\end{equation}
\begin{itemize}
  \item \emph{Serie impar.} $a_3=0\cdot a_1=0$, y con él todos los impares
        posteriores: la serie impar se reduce a $y_2=x=P_1(x)$.
  \item \emph{Serie par.} $a_2=-a_0$, $a_4=\tfrac13a_2=-\tfrac13a_0$,
        $a_6=\tfrac35a_4=-\tfrac15a_0$ y, por inducción,
        $a_{2k}=-a_0/(2k-1)$ para $k\ge1$. Por tanto
        \begin{equation}\label{eq:legendre-q1}
          y_1=1-\sum_{k=1}^{\infty}\frac{x^{2k}}{2k-1}
             =1-x\,\artanh x
             =1-\frac{x}{2}\,\ln\frac{1+x}{1-x},\qquad\abs{x}<1 .
        \end{equation}
\end{itemize}
La función $y_1$ (que es $-Q_1$, con $Q_1$ la función de Legendre de segunda
especie) diverge logarítmicamente en $x=\pm1$: la cota del \cref{teo:fuchs}
es \emph{óptima} en este caso. Obsérvese también que la ecuación
$(1-x^{2})\,y''-2x\,y'+\tfrac34y=0$ del \cref{ej:ordinario} es la de Legendre
con $\ell(\ell+1)=\tfrac34$, es decir, $\ell=\tfrac12\notin\N$: ninguna de sus
dos series se corta y ambas tienen radio exactamente $1$.
\end{ejemplo}

\begin{observacion}[Origen físico y «cuantización»]\label{obs:legendre-fisica}
La ecuación de Legendre aparece al separar variables en la ecuación de
Laplace en coordenadas esféricas, con $x=\cos\theta$. Los puntos $x=\pm1$
corresponden a los polos $\theta=0,\pi$, donde la solución física debe ser
finita. Por \eqref{eq:legendre-q1} y la \cref{prop:polinomios-legendre}, eso
sólo es posible si una serie se trunca, lo que exige $\ell\in\N$: la
condición de regularidad \emph{cuantiza} el parámetro, exactamente como la
condición de normalizabilidad cuantiza la energía en la ecuación de Hermite
del oscilador cuántico \cite{TEMAV}, \cite[§13.1]{ARFKEN}.
\end{observacion}
